\section{Research gaps}
\label{sec:gaps}
Despite the rapid development of cloud computing technologies and the increasing adoption of DevSecOps practices, existing research has primarily focused on individual technologies, security mechanisms, or specific stages of the software development lifecycle. Multiple literature reviews consistently report that the current body of research remains fragmented across several dimensions, including security automation, tool integration, operational workflows, and empirical validation. Myrbakken and Colomo-Palacios observed that DevSecOps research predominantly concentrates on practices and tools rather than comprehensive operational frameworks \cite{Myrbakken2017DevSecOps}. Similarly, Rajapakse et al. identified continuous security integration, automation, and security orchestration as persistent research challenges through a systematic review of DevSecOps studies \cite{Rajapakse2021DevSecOpsSLR}. More recently, Zhao et al. concluded that opportunities remain for broader operational integration and more comprehensive validation of DevSecOps approaches across complex computing environments \cite{Zhao2024DevSecOpsMLR}. These observations suggest that several important research gaps remain, particularly for security operations within hybrid cloud environments.

\vspace{0.3cm}

\textit{Gap 1: Limited Operational Frameworks for Hybrid Cloud Security}.
Existing research has proposed numerous DevSecOps practices, cloud security frameworks, and automated security solutions to improve secure software delivery. However, these studies primarily focus on software development activities, security tools, or individual operational practices instead of providing an integrated operational process capable of coordinating infrastructure management, software deployment, and continuous security within hybrid cloud environments. While the NIST Cloud Computing Reference Architecture defines the architectural relationships among cloud components and describes the characteristics of hybrid cloud infrastructures, it does not prescribe an operational process for integrating system administration, software engineering, and security activities \cite{Liu2011NISTRA}. Likewise, both Myrbakken and Colomo-Palacios and Zhao et al. emphasized that existing DevSecOps research is largely organized around practices, tools, and lifecycle activities, whereas comprehensive operational frameworks remain comparatively limited \cite{Myrbakken2017DevSecOps,Zhao2024DevSecOpsMLR}. This indicates a need for research that investigates integrated operational processes tailored specifically to hybrid cloud environments.

\vspace{0.3cm}

\textit{Gap 2: Insufficient Integration of Continuous Security Throughout Hybrid Cloud Operations}.
Continuous security has become a fundamental principle of secure software engineering, and the NIST Secure Software Development Framework recommends integrating security practices throughout software development and deployment activities to improve software integrity and reduce vulnerabilities \cite{NISTSSDF2022}. Nevertheless, implementing continuous security remains a persistent challenge in practice. Rajapakse et al. identified automation of security activities, continuous security assessment, and balancing software delivery speed with security assurance as recurring challenges across existing DevSecOps studies \cite{Rajapakse2021DevSecOpsSLR}. Similarly, Myrbakken and Colomo-Palacios reported that security integration across the software lifecycle remains one of the central themes requiring further investigation \cite{Myrbakken2017DevSecOps}. While current research predominantly concentrates on integrating security into CI/CD pipelines, comparatively fewer studies investigate how continuous security should be coordinated across both software engineering activities and infrastructure operations in heterogeneous hybrid cloud environments. Consequently, further research is required to establish operational processes that embed continuous security throughout the entire hybrid cloud operational lifecycle.

\vspace{0.3cm}

\textit{Gap 3: Limited Orchestration of Heterogeneous Security Tools and Security Findings}.
Modern software projects increasingly employ multiple security assessment tools, including static application security testing, software composition analysis, dependency vulnerability scanning, and infrastructure security analysis. Although these tools improve vulnerability detection from different perspectives, they typically operate independently and generate heterogeneous outputs with different reporting formats, severity classifications, and confidence levels. Zhao et al. observed that existing DevSecOps literature primarily categorizes research according to security practices, tools, and metrics, while comparatively less attention has been devoted to the orchestration and integration of heterogeneous security tools into unified operational workflows \cite{Zhao2024DevSecOpsMLR}. Rajapakse et al. likewise identified tool integration and security automation as recurring research challenges affecting DevSecOps adoption \cite{Rajapakse2021DevSecOpsSLR}. These findings suggest that systematic approaches for coordinating multiple security tools, consolidating their outputs, and supporting unified operational decision-making remain relatively underexplored.

\vspace{0.3cm}

\textit{Gap 4: Limited Empirical Validation of Integrated Security Operational Processes}.
Existing literature has proposed numerous recommendations, conceptual architectures, and best practices for implementing DevSecOps. However, comprehensive empirical validation of integrated operational processes remains comparatively limited. Myrbakken and Colomo-Palacios highlighted that much of the current literature discusses DevSecOps from conceptual or practical perspectives without providing comprehensive evaluation of integrated operational frameworks \cite{Myrbakken2017DevSecOps}. More recently, Zhao et al. identified empirical validation as one of the research dimensions requiring further attention, particularly for studies involving broader operational integration and practical implementation \cite{Zhao2024DevSecOpsMLR}. Furthermore, Rajapakse et al. emphasized the need for additional empirical studies to evaluate integrated DevSecOps approaches in realistic operational environments \cite{Rajapakse2021DevSecOpsSLR}. Therefore, there remains a need for practical implementations that demonstrate how integrated security operational processes can be deployed and evaluated within hybrid cloud environments.